\chapter{引言}

\section{研究背景}

作为大学里数学专业学生的专业教材，《数学分析》华东师范大学编（第四版）主要用于数学分析课程的学习。17世纪时，牛顿和莱布尼兹等人开创了数学分析，到了19世纪则是由魏尔斯特拉斯和柯西等人对此进行奠基。

数学分析是以函数为基础，从局部和整体两个角度对函数的基本形式进行研究，并由此构成积分和微分的基础。微分研究函数的局部性质，如变化率；研究范围包括微分的性质、计算和直接应用；基本概念为导数和微分，微分法即是求导的过程。区别于微分，积分立足于整体上微小变化积累的程度与效果，尤侧重非均匀的变化的研究；研究范围包括积分的性质、计算和直接应用；基本概念为原函数和定积分，积分法即是求积分的过程。谈及历史，牛顿、莱布尼兹等数学家，在1670年前后，创立了一套关于求导数与积分的基本定律，并提出了求导数与积分的互逆运算，创造出了牛顿·莱布尼兹公式对上述定律与运算进行表示，从而合并了微分与积分，成立了一个新的分支—微积分学。而后，又在他们的后继数学家的推动下，如欧拉，让原本只有极少数学家能理解的一系列简单而具体的问题，变成了一种普通人只要稍加练习就能学会的方法，从而在各种科学和技术上获得极大的发展。因此，微积分的出现与发展被认为是人类文明史上划时代的数学事件之一。数学分析中，无穷级数也有所涉及。历史上，无穷级数的使用由来已久，但只在成为数学分析的一部分后，才得到真正的发展和广泛应用。在数学分析中，无穷级数与微积分从来都是密不可分和相辅相成的，但与积分所具备的连续性相比，无穷级数虽也是微小量的叠加和积累，采取了离散的形式。

《数学分析》华东师范大学编（第四版）共二十三章，分别置于上下两册。与上文类似，全书以实数集与函数开头，对数列以及函数的极限进行定义探寻，导数和微分与高中知识接轨，接着从不定积分、定积分、反常积分中对积分进行系统研究。紧接着级数、多元函数的极限与连续性、多元函数微分学登场。而后含参量的积分、重积分是积分之于几何的联系，曲线积分加深了此般关系,曲面积分又是线到面的推广。数学分析课程的主要内容便由上述内容构成。全书奠定了后续更为细化且深入的数学分支的学习的基础，知识丰富，结构紧凑。同时，该书的内容设置，在一定程度上也展现出数学分析这门学问变迁的脉络。



\section{研究发展现状}


不等式有关理论作为古老的自然科学一大分支，数年来都备受学者青睐。总的而言，世界各国对于不等式的研究越来越高深，越来越复杂与多维。纵观古今，不等式领域中少不了中国数学家的身影，如华罗庚先生；近年来，在国际上，不等式和其应用的领域中也有不少中国学者的贡献。单一类别的不等式研究，如Hilbert不等式，对其参量化、积分形式、积分形式的推广等都有不同研究报告的呈现。再到各种不等式的证明方法的多样化研究，如针对中等数学的用向量法、概率方法等证明不等式。而后高等数学中运用数学分析的方法进行证明。再有将数学文化融入不等式的研究与教学中，如从信息技术入手和课堂艺术出发……不等式的研究多维且深入，拓展到生活的方方面面。

均值不等式的研究发展现状可谓是极为全面，其妙用不只是在数学方面呈现，物理等其他学科也均有涉足。均值不等式的研究涵盖了其适用条件、性质、证明方法等诸多方面。均值不等式也被运用到求解不同类型的题目中，如最值问题、数列问题等\cite{moslehian2021arithmetic}。戴伟在《关于均值不等式的一个证明》中更是利用均值不等式给出了酉矩阵的一个刻画\cite{戴伟2020关于均值不等式的一个证明}。谈及国外，Davit Harutyunyan加强了几何-算术平均值不等式\cite{harutyunyan2018cauchy}。Mohammad Sal Moslehian等人更是给出了不定型的算术-几何均值不等式\cite{moslehian2021arithmetic}。

伯努利不等式由瑞士数学家伯努利提出，其一般形式笔者将在后文中给出。国外的研究中，其涉及到原函数和双曲函数等多种函数形式，也对负指数形式进行了探索。在国内，马玉梅等在《微积分教学中几个问题的思考》将伯努利不等式应用到Jacobsthal不等式\cite{马玉梅2020微积分教学中几个问题的思考}。该篇更是对 $n$ 的奇偶性进行了探讨。杨克昌在《一类不等式的加强与综合》中利用伯努利不等式解决含参变量的不等式\cite{杨克昌1997一类不等式的加强与综合}……伯努利不等式的研究主要运用到其推广形式，在高深的数学领域中依旧有深入探讨的必要。

柯西不等式的一般形式及变式形式涉及多种，后文中笔者会给出一二。柯西不等式就其发展进程、证明方法、运用形式、改进方式都是研究的方向。陈明在《柯西不等式的几点注记》中展现出柯西不等式在数学理论中的重要地位\cite{陈明2018柯西不等式的几点注记}。柯西不等式的证明，以及其在函数求最值、证明不等式、几何上的广泛应用在也本篇中得到了阐述。
千变万化的多项式也是在不等式的不断发展中逐步步入人们的学习与生活中的。均值不等式、伯努利不等式、柯西不等式作为最基础的三类重要的不等式，融入到各类不等式的研究中，体现出不同性质下的价值。

\section{研究内容及方法}

本文主要对《数学分析》华东师范大学编（第四版）中所涉及的三类不等式：均值不等式、伯努利不等式、柯西不等式的证明、推广变式及例题进行分析，并进一步给出教材以外的相应高中数学或大学数学例题，展现出不等式在不同阶段的运用，展示出不等式的教育价值，使得学生更能掌握对于这三类不等式的理解与应用。

通过对以上所涉及知识的研究与分析，本文运用的主要方法为文献参阅法：通过查阅资料研究上文所提及的三类不等式，了解其发展，掌握其证明，熟悉其应用，查阅的主要资料为期刊论文、博硕士学位论文等。



\section{研究目的及意义}

不等式具有研究的重大价值，这不仅是从基础性出发，也是于重要性而言。在解题中巧妙运用不等式总会让人眼前一亮，但碍于自身数学能力的不足，总会有很多人难以发现不同不等式的构造与使用方法，尤其是当题目变得复杂时。本文具体从《数学分析》华东师范大学编（第四版）出发，总结出三类不等式:均值不等式、伯努利不等式、柯西不等式，将其在书中例题的应用整合并推广，体现出该三类不等式在数分中的重要性的同时，给出更多相应例题，让例题从高中数学跨越至大学数学，不仅展现出三类不等式在数学分析教材外的突出作用和应用办法，也表现出三类不等式在不同阶段下的不同妙用，从而整体达到对不等式及其应用更深入的认识，体现出不等式的应用和教育价值。